\chapter{Implementation}
\label{ch:Implementation}

\label{sec:ImplementationTarget}
This research aims to scale ORAMs, using \textit{3-Party DUORAM} as an example, arbitrarily.
The already existing \textit{DUORAM} implementation by Vadapalli et al. \cite{vadapalli2023duoram} is optimized for 64-bit values.
Larger values are not supported yet.
After explaining the existing code, this chapter highlights and discusses the most promising approaches.

\section{Scaling Based on the Protocol}
\subsection{Scaling in Theory} \label{sub:ScalingInTheory}
This section conducts a theoretical analysis of how to arbitrarily scale \textit{DUORAM} to fulfil the objective of this research.
It is based on reading and writing protocol from \ref{sub:AdvancedReadingAndWriting} taken from Vadapalli et al. \cite{vadapalli2023duoram}.

\subsubsection*{General Assumptions} \label{subsub:ScalingInTheoryGeneralAssumptions}
This paragraph highlights general thoughts and assumptions on arbitrarily scaling duoram entries.
The affected values are independent of the phase and operation executed, thus placed separately in this subsection.

The theoretical protocol is not bound to computational limits.
Therefore, the database $D$ is unlimited in the number of entries it can hold, solely the entry size is fixed.
This means, scaling the database equals scaling the entry size of every entry $D[i]$.

As both $D$ and its blinded equivalent $\widetilde{D}$ with its respective additive shares $\widetilde{D}_0,\widetilde{D}_1$ are of the same size, the blinded shares will automatically adjust their size once $D$ is scaled.
This will also affect its respective shares $D_0,D_1$ which are the same size as $D$.
Since $\widetilde{D}_b\leftarrow D_b+\xi_b\text{ for } b\in\{0,1\}$ the blinding factors $\xi_0,\xi_1$ will grow in their entry size too.
Note that all instances of $D,D_0,D_1,\xi_0,\xi_1,\widetilde{D}_0,\widetilde{D}_1$ are of equal size.
They therefore grow at the same rate when the size of $D$ is adjusted keeping the same size.

These general assumptions are taken as granted in this theoretic part.

\subsubsection*{Scaling during Preprocessing} \label{subsub:ScalingDuringPreprocessing}
As preprocessing for reading and writing is nearly similar in their goals and processes, both preprocessing phases are covered in this subsection.
As the main goal of preprocessing is to generate RDPF-Triples for each peer, this process is affected the most during scaling and will affect its derived values.

\paragraph{Scaling DPF Generation}\label{par:ScalingInTheoryDPFGeneration}
The theoretic DPF generation process follows the introduced construction of Boyle et al. \cite{boyle2016function} from \ref{sub:BGIDPFConstruction}.
It takes an $n$-bit seed as input and recursively uses a length-doubling PRG to generate two binary trees.
With correction words, flag bits, and tree invariants both trees are adjusted to represent a DPF.

Vadapalli et al. \cite{vadapalli2023duoram} extend this concept by further using the results of this DPF construction.
Each peer is assigned three DPF shares in form of RDPF-Triples.
These are then evaluated leaving each peer with three vector pairs $(v_b^k, t_b^k)$ with the known properties.

To reconstruct the effects of adjusting the database entry-size the author chose to traverse the algorithm of the BGI-construction (\ref{sub:BGIDPFConstruction}) in reverse.
Recall the properties of $v_b^k$ representing $v^k = e_{i^*}\cdot M$ when combined.
This means $v^k$ is a vector holding $M$ at the target index $i^*$ and zero everywhere else.
When additively shared among both peers, it is divided into two vectors $v_0^k,v_1^k$ holding additive shares of $M$ at the target index and random values everywhere else.
These random values will add to $0$ once combined.
As each entry has the same size and $v_b^k[i^*]$ has size $|M|$ each entry of $v_b^k$ has the same size as the message.
As $D[i]$ is adjusted in its size, $M$ will adjust in the same manner allowing the client to write messages of the new size.
This means each entry $v_b^k[i]$, having the same size as $M$, will also adjust its size once $D[i]$ has changed.

Recall $v_b^k[i]$ being the leaf label at index $i$ from $P_b$'s $k$'th DPF-tree.
As $v_b^k$ adjusts to the size of $M$, each leaf label will have the size of $M$.
Since the BGI-construction uses a length-doubling PRG, the seed of the tree will have size $|M|$ too.
A single execution of the length-doubling PRG therefore generates two children of size $|M|$ each.
This means a single length-doubling PRG execution on a seed of size $|M|$ results in a string of size $2|M|$

As mentioned previously, $F^k_b$ is the negative sum of the according leaf labels $-\sum^n_i{v_b^k[i]}$.
Thus, $F$ will also adjust once $D[i]$ has adjusted its size.

$t_b^k$ represents the flag bits of each leaf in $P_b$'s $k$'th DPF-tree.
As each leaf holds a single flag bit $t_b^k$'s size is solely dependent from the amount of leafs in the DPF.
This means the size of $t_b^k$ will not be affected by scaling $D[i]$.

To conclude, adjusting the entry size of $D[i]$ results in $v_b^k,F^k_b,\text{ and }M_b$ adjusting their size in the same manner.
Changing the size of these values will result in size-changes while generating DPFs and in the preprocessing phase in general.

\subsubsection*{Online-Phase Scaling} \label{subsub:ScalingDuringOnlinePhase}
The online-phase does not yield any new values but the indices $ri$ and $i$.
They indicate which element of the database to retrieve, thus affecting the offset $S$ during cyclic shifting.
This, however, is not dependent from the size of entry $D[b]$ but the amount of entries.
If the datatype for the indices $i$ and $ri$ is large enough to uniquely identify each leaf, no changes are necessary. 

\section{Existing Implementation}
\label{sec:ExistingImplementation}

\section{Preliminaries}\label{sub:ImplementationPreliminaries}
\subsubsection*{Representing DPFs} \label{subsub:implementationRepresentingDPFs}
Both, Distributed Point Functions (\ref{sub:DistributedPointFunctions}) and Distributed Comparison Functions (\ref{sub:DistributedComparisonFunctions}), inherit from the same parent class \code{dpf}.
Again, recall the representation of point- and comparison-functions as binary trees.
Thus, the most trivial representation of a single point function is saving the seed of the root node and generating the whole tree every time the point function is used.
Since this approach is too computational expensive for repetitive evaluations, DPFs are stored in an already evaluated form with the possibility to recreate the tree when needed.

Therefore, the base class \code{dpf} stores the seed of the root node.
Also, it stores a vector \code{cfbits} which holds all flag bits on the path to target index $\alpha$.
Bit $i$ of this vector represents the bit of the $i$'th node on the path to $\alpha$.
It also implements the functionality to descend a tree given a seed and a bool indicating which child to calculate.
The code further defines special DPFs namely \textit{Distributed Point Functions for Memory Access (RDPF)} and \textit{Distributed Comparison Functions (CDPF)}.

\paragraph{RDPFs}\label{par:preliminariesRDPFs}
RDPFs inherit from the basic DPF implementation.
It expands its parent with further attributes and the \code{LeafInfo} struct.
\code{leaf\_cfbits} stores the correction-flag bits for the leaf level.
Recall the correction-flag bits being the low words of the labels (\ref{sub:BGIDPFConstruction}).
This is equivalent to $t_b$ from the high-level \textit{DUORAM} overview in \ref{sec:3pDuoram}, which means bit $i$ holds the correction flag-bit of leaf $i$.
\code{li} is a vector holding \code{LeafInfo} structs.
In normal RDPFs, which is the standard in this research, it will only hold a single \code{LeafInfo} struct.
\code{LeafInfo} is essential for using the evaluated RDPF result in a read or write operation.
It holds the correction word for the leaf level, which itself is a XOR-share, and all further attributes to easily convert it into an additive share.
This is essential, since \code{duoram} is based on additive shares. 
Thus converting from XOR RDPF-shares to additive shares simplifies further evaluation processes a lot.
More information on \code{LeafInfo} and the concept of \code{DPFNode}s and \code{LeafNode}s is given in \ref{subsub:DpfNodeAndLeafLabel}.

The code stores RDPFs as triples, called \code{RDPFTriple}.
As the name suggests, it holds three RDPFs, each with the same target, again, either XOR or additively shared.
To generate a RDPF-Triple including three RDPFs on a given target, one uses its constructor which itself calls the RDPF's constructor three times.
All three RDPFs are stored in a vector, the \code{RDPFTriple}.
The triple itself is stored in the \code{mpcio} class assigned to each player.
When generating these triples in the preprocessing phase, they are written to the background memory and called automatically when performing reads or updates in the online phase.
This construction makes it easy to keep track of the used RDPFs.

\paragraph*{Leaf Construction}\label{par:leafConstruction}
Leafs are stored as concatenations of 64-bit vectors.
A single leaf holds the functionality for both reading and writing.
This means it holds 64-bit for writing, which will add up to $M_b$ at the target index, and 64-bit for reading, adding up to the shifting vector when combined with the other peer's shifting vector.
Concatenating both 64-bit values in a single value results in a single 128-bit vector.
Writing bits are placed at the upper 64-bits, the lower word represents shares of the shifting vector.
Therefore, when executing a read and requiring the flag bits, one can simply add up (or XOR) the least significant bits of all the leafs without performing additional shifting operations.
This means a tree of depth $d$ holds $2^d$ leafs, each of size 128.

\paragraph{CDPFs}\label{par:preliminariesCDPFs}
CDPFs too, inherit from the base structures predefined by \code{dpf}.
They are, however, stored as tuples.
One can choose whether to share the target additively via \code{as\_target} or distribute XOR-shares, called \code{xs\_target}.
This process is similar to RDPFs, since both CDPFs within a tuple hold the same target.
Thus, combining the targets of $P_0,P_1$ results in the actual target index.
A tuple of CDPFs allows the client to compare a given value to $0$ as described in \ref{sub:CDPFConstruction}.

\subsubsection*{types.hpp} \label{subsub:implementationTypeshpp}
This file defines basic macros, aliases, structs, and their operators.

The most prominent type-aliases are:
\begin{itemize}\label{itemize:importantAliases}
    \item \code{nbits\_t}: 
    This value counts the number of bits in a value, which means \code{value\_t}.
    Thus, the chosen datatype needs to be large enough to store \textit{VALUE\_BITS} bits.
    The default value is \code{uint8\_t}.
    \item \code{bit\_t}:
    Chooses how a single bit is represented, usually in the \code{RegBS} struct.
    The default datatype for this alias is \code{bool}.
    \item \code{address\_t}: 
    Decides the size of the addresses in the secret-shared memory.
    Is set to \code{uint32\_t}.
    \item \code{DPFnode}: 
    Implements the datatype of a DPFnode and establishes the basis for defining Leaf nodes.
    It is of type \code{\_\_m128i}.
    Leafs are also implemented using this type.
    Evaluation on leaf nodes splits the vector in a lower and upper half which makes it two 64-bit vectors.
    More information on this alias is provided in \ref{subsub:DpfNodeAndLeafLabel}.
    \item \code{value\_t}: 
    Decides the datatype of values in the secret-shared memory.
    The default value is \code{uint64\_t}.
    \item \code{AES\_KEY}:
    Represents an AES-key for 128-bit AES.
    More details on this alias are found in \ref{subsub:implementationAesni} and \ref{sub:AESSizeLimitations}.
\end{itemize}
Important macros are:
\begin{itemize}\label{itemize:importantMacros}
    \item {VALUE\_BITS}: 
    This will decide the size of \code{value\_t}. 
    The default value is 64.
    \item {ADDRESS\_MAX\_BITS}: 
    Decides the size of \code{address\_t}.
    The default value is 32.
\end{itemize}
The most prominent structs are:
\begin{itemize}\label{itemize:importantStructs}
    \item \code{RegAS}: 
    Represents a register holding an additive share of type \code{value\_t}.
    \item \code{RegXS}: 
    Represents a register holding an XOR share of type \code{value\_t}.
    \item \code{RegBS}: 
    Represents a register holding a bit-share of type \code{bit\_t} which itself is a \code{bool} type-alias.
\end{itemize}

\subsubsection*{Intel AES-NI} \label{subsub:implementationAesni}
The existing source code is optimized for Intel's \textit{Advanced Encryption Standard New Instruction Set (AES-NI)}\footnote{Intel's publication on AES-NI is published on their \hyperlink{https://www.intel.com/content/dam/doc/white-paper/advanced-encryption-standard-new-instructions-set-paper.pdf}{website} (https://www.intel.com/content/dam/doc/white-paper/advanced-encryption-standard-new-instructions-set-paper.pdf)}.
As the name suggests, it extends and optimizes Intel's existing instruction set for AES-operations.
These improvements come in form of the \code{wmmintrin.h} and \code{emmintrin.h} header files implementing core functions of \textit{AES-ECB (Advanced Encryption Standard in Electronic Code Book Mode)} based on \textit{AES-NI}.
As this header file is included in the gcc compiler, no additional installations for using both headers are required.
Yet, the flags \code{-maes}, and \code{-msse2} need to be included for the linker to work.
Since AES-ECB is of great use when generating DPFs, these optimizations will foster great performance improvements.
Operations like \code{AES\_key\_gen\_assist} are optimized for the \code{\_\_m128i} \textit{Single Instruction Multiple Data (SIMD)} datatype, representing a 128-bit vector.
Therefore, the implementation stores any data used in AES-contexts in this format.

The careful reader notices \code{\_\_m128i} vectors being twice as large as the \code{uint64\_t} type-alias \code{value\_t} (\ref{itemize:importantAliases}).
This is due to AES-ECB being used as a \textit{Pseudo Random Generator (PRG)} during creation and traversion of DPF trees.
The code uses its security properties of being computationally random \cite{rijmen2001advanced} to calculate new labels given a seed and key.
Both parties use the same private key, thus AES-ECB becomes pseudorandom.
Labels are stored as \code{\_\_m128i}, leafs as explained in the previous section \ref{par:leafConstruction}.
Therefore, by using a \code{\_\_m128i} seed and labels encrypting each label using AES-ECB recursively, one simulates a length doubling pseudo random generator \cite{vadapalli2023duoram} as required to run the BGI-DPF-Construction (\ref{sub:BGIDPFConstruction}).
More on label and leaf construction can be found in the following sections.

AES for 128-bit uses ten encryption-rounds \cite{rijmen2001advanced}.
Therefore, after inserting the seed which resembles the key of the AES operation, ten more sub-keys are generated starting from the seed.
This structure is stored in format of an \code{\_\_m128i[11]} array.
The seed at index $0$, the ten sub-keys at the following indices.
It is represented as type-alias \code{AES\_KEY} which is equivalent to \code{\_\_m128i[11]} (\ref{itemize:importantAliases}).
To simplify key-creation, \code{AES\_KEY} is wrapped in the struct \code{PRGKey}.
It takes a single \code{\_\_m128i} value as constructor input and generates ten more sub-keys using the input as seed for AES-ECB.
As the input and the key are set to the same value every time, one receives a pseudorandom number generator.

For AES to be pseudorandom, keys need to be the same throughout instances.
Thus, all instances use the same \code{PRGKey} instances for creating labels and leafs, called \code{prgkey} and \code{leafprgkeys[3]}.
Note \code{leafprgkeys[3]} being an array holding three \code{PRGKey} structs.
This distinction between keys for label and leaf generation is necessary as one needs to differentiate cases handling different values for \code{LWIDTH} (\ref{subsub:WIDTHandLWIDTH}).
A non-standard value of \code{LWIDTH}, which means any value $x\neq1$, will require multiple keys in the leaf layer as several Leaf labels need to be instantiated (\ref{subsub:DpfNodeAndLeafLabel}).
As this matter is only present in the leaf layer, using separate keys during generation will increase performance and readability of the code.
The implementation uses the first 20 digits of $\pi$ for \code{prgkey} and 20 digits of $e$ for each instance of \code{leafprgkeys}.

Requirements derived from this, along with further details on \code{WIDTH} and \code{LWIDTH} are described in \ref{sec:GeneralLimitations}.

One practical drawback of using \code{\_\_m128i} vectors is its debugging incompatibility when using \textit{Valgrind}\footnote{More information can be found on their website and documentation at https://valgrind.org/}.
\textit{Valgrind} is a debugging suite designed for spotting memory leakages in the code.
It is, however, not compatible for using SIMD datatypes which triggers an error and ultimately stops program execution.

\subsubsection*{Duoram and Shapes} \label{subsub:implementationDuoram}
Each party is assigned a \code{duoram} object.
It holds the following information: 
\begin{itemize}
    \item player: 
    The number of the party.
    0 or 1 for computational peers and 2 for the server.
    \item database: 
    Vector of registers representing the party's share of the database $D_b$.
    These entries are either additively (\code{RegAS}) or XOR shared (\code{RegXS}).
    Entries are stored as \code{RegAS} by default and can be changed manually.
    This means the database is a vector holding $n$ register entries.
    \item blind: 
    The blinding factor $\zeta_b$ to blind its own database.
    As the blinding factor is continuous, it has the same size as the database.
    \item {peer\_blinded\_db}:
    This is the blinded database of party $P_b\text{ with }b\in\{0,1\}$ the other party $P_{b'}\text{ with }{b'-1}\in\{0,1\}$.
    \item Party $P_2$ uses the attributes \code{p0\_blind}, \code{p1\_blind} to represent the blinded databases of both parties $P_0,P_1$.
    The size is, again, the same as the database.
\end{itemize}

Each \code{duoram} must be accessed via a hierarchy of \textit{Shapes}.
The implementation uses a hierarchy consisting of solely a single shape called \code{Flat}.
A shape provides basic functions for reading and updating the duoram.
This is primarily the index operator \code{[]} for reading and the \code{+=} operator for updating.
This has the great advantage of updating accessing contexts while keeping the underlying database consistent.
For example, when executing multiple database accesses multiple RDPF-Triples and thus separate contexts are needed as keeping the context might leak accessing-information due to reusing the same RDPF-Triple.
Instead of assigning these contexts to the same database introducing a great overhead one simply updates the underlying shape introducing a fresh context.
Without a shape, one manually has to update the database, reassign new RDPF-Triples to each player and manipulate their \code{mpcio} objects.
\code{context()} as defined in \code{Flat} generates a new shape including fresh RDPF-Triples and assigns it to each party.
This reduces the overhead to a minimum and introduces independent reads and writes.

The base implementation uses the \code{Flat} shape only.
Extensions on different shapes are not planned or necessary, thus no further descriptions on different shapes and shape hierarchies are required.

\subsubsection*{WIDTH and LWIDTH}\label{subsub:WIDTHandLWIDTH}
RDPFs can be reused for reading.
This, however, will leak the information that reads were executed on the same location multiple times.
The location stays secret \cite{vadapalli2023duoram}.

This process is similar for writing, yet it requires a few code adjustments.
In short, one simply extends the given leaf construction (\ref{par:leafConstruction}) by adding as many 64-bit values representing shares of $M$ as wanted.
Vadapalli et al. chose to template their \code{RDPF} struct with a property called \code{WIDTH} of type \code{nbits\_t} (\ref{itemize:importantAliases}).
\code{WIDTH} is set to $1$ by default.
This indicates one write access per RDPF, no reuse is possible.
Increasing this value to $n$ adjusts the type aliases \code{RegASW} and \code{RegXSW}, both containing $n$ instances of \code{RegAS} and \code{RegXS} respectively.
On a high level, this means that each RDPF can hold up to $n$ different values for the message $M$.
Therefore, by increasing \code{WIDTH} to a value of $n$, one can reuse a given RDPF for up to $n$ different writing accesses.

\code{LWIDTH} is directly derived from \code{WIDTH}.
It is calculated by $1+\text{WIDTH}/2$.
As it is of type \code{nbits\_t} (\ref{itemize:importantAliases}) setting \code{WIDTH} to $1$ will set \code{LWIDTH} to $1$ too.
Recall leaf nodes being stored as an array of \code{\_\_m128i} values, with 64-bit reserved for reading and 64-bit reserved for writing.
Performing either a single read or write will only require the array to contain only a single \code{\_\_m128i} value.
Increasing \code{WIDTH} to $n$ will lead to $n$ 64-bit values being needed for $n$ writes.
Therefore, the array representing leaf nodes must be adjusted to hold $n$ \code{\_\_m128i} values.
\code{LWIDTH} represents the number of \code{\_\_m128i} values necessary to accommodate \code{WIDTH} many writes per instance.

Limitations on these values are listed in \ref{sub:LimitWIDTHandLWIDTH}, more information on this construction is given in the following section.

\subsubsection*{DPFNode and LeafNode}\label{subsub:DpfNodeAndLeafLabel}
The concepts of DPF nodes, leafs, \code{WIDTH}, and \code{\_\_m128i} vectors are closely intertwined.

As already mentioned, DPF trees are generated using AES-ECB (\ref{subsub:implementationAesni}).
For the goal of storing 64-bit values in the base implementation, one needs to simulate a length-doubling PRG.
This can be done by using AES-ECB on vectors twice as large as the actual label, which means 128-bits.
Therefore, a single node for a particular DPF tree storing values of size 64 has nodes of size 128.
Vadapalli et al. consequently drew the logical conclusion to use \code{DPFnode} as type alias for \code{\_\_m128i}.
As Intel's AES-ECB implementation (\ref{subsub:implementationAesni}) directly uses \code{\_\_m128i} vectors as input and output, \code{DPFnode} can be used for AES operations without conversion.

Recall the definition and usage of \code{WIDTH} and \code{LWIDTH} from \ref{subsub:WIDTHandLWIDTH}.
In short, each leaf can be divided into $2n*$\code{LWDITH} blocks of size 64, which makes it reusable for \code{WIDTH} many updates.
This means, to implement a DPF tree with a certain value of \code{WIDTH}, one needs a leaf label of size $\code{LWIDTH}*128$.
Thus, a \code{LeafNode} is chosen to be a type alias for a vector holding \code{LWIDTH} many \code{\_\_m128i} values or \code{DPFnode}s respectively.
This results in the type definition \code{using LeafNode = std::vector<DPFNode, LWIDTH>}. 

\subsubsection*{mpcio and Communication} \label{subsub:implementationMpcioAndCommunication}
All three parties execute the same executable \code{prac} which can be generated by using the \code{MAKEFILE}.
These executables communicate asynchronously via TCP-sockets.
When executed locally, which this research assumes to be the default, the peers addresses are all hosted on the \code{localhost}.
Sockets, player information, sent and received messages and the functionality to send and receive are implemented in the \code{MPCPeerIO} struct.
\code{MPCPeerIO} can be freely translated to \textit{"Input and Output Functionality for performing MCP communication with computational peers"}.
It is one of the first structs to instantiate and will be used to create most of the other objects.

\code{MPCPeerIO} is the child of the \code{mpcio} struct which holds general io-attributes and also provides functionality for timing the communication process and timing of round-trip-trimes.
This struct also holds information about the chosen processing mode.
The program offers three possibilities for this value: 
\begin{itemize}
    \item -p: 
    Defines the \code{PRECOMPUTATION\_MODE}, which will execute the preprocessing phase as explained in section \ref{subsub:DuoramPreprocessing}.
    \item leaving the specifier blank: 
    Will execute the \code{ONLINE\_MODE}.
    Note that this requires the preprocessing mode to be executed first since it requires dpf-triples for reads and writes.
    Thes user can specify which tests to execute.
    The default test is called \code{duoram} and will check reads and writes, each independant and interleaved.
    Read more about testing in section \ref{sub:testingAndBenchmarking}.
    \item -o:
    Will execute the \code{ONLINE\_ONLY} mode.
    This means no precomputation will take place, all needed resources are computated on the fly.
    Of course, this will impact performance but will come in handy for performing smaller tests.
    With this specifier, too, one can choose which tests to execute.
    These tests will be the same as in the online mode.
\end{itemize}

All three parties will use a predefined message structure.
Therefore, generated precomputation \ref{subsub:implementationPreprocessing} data will implement this protocol.
It separates a message into byte-segments with each segment accomodating a special type of data:
\begin{itemize}\label{itemize:preprocessingStructure}
    \item 0x01: Indicates the type.
    \item 0x01 to 0x30: RDPFs. Since the size of RDPFs depends on its depth the given space is defined in this manner.
    \item 0x40: CDPFs 
    \item 0x80: Multiplication Triple for performing MPC-multiplications (\ref{sub:MultiplicationAdditionTriples})
    \item 0x81: AND Triple for performing MPC-AND-operations (\ref{sub:MultiplicationAdditionTriples})
    \item 0x83: Select Triple. Since it will not be uesd further the benevolent reader may read the comments in \code{types.hpp}.
    \item 0x8e: Counter counting the AES-operations. This can be used for testing.
    \item 0x00: Shows the end of preprocessing. 
    \item One extra byte is reserved for \code{subtype}.
    More details about this byte can be found in \ref{sub:LimitWIDTHandLWIDTH}.
\end{itemize}

More detailed information about the TCP-sockets and the underlying functions is given in \ref{subsub:gmpCommunication}.
\section{Scenarios and Execution} \label{sub:implementationScenarios}
As specified in the previous paragraph \ref{subsub:implementationMpcioAndCommunication}, \code{prac} is executed three times simultaneously.
Hence, assuming one executes locally without using docker, one needs to run \code{prac} in three terminals at once.
Calling \code{prac} requires several arguments: player number, which mode to execute (-p, -o, or blank), number of threads for faster execution (-t), and for the server $P_2$ the type, amount and depth of functions to calculate.
\code{r} will represent RDPFs (\ref{sub:RDPFTriples}).

Thus, the function calls for preprocessing will look as follows:
\begin{enumerate}
    \item \code{./prac -t 3 -p 0}
    \item \code{./prac -t 3 -p 1 localhost}
    \item \code{./prac -t 3 -p 2 localhost localhost r20:20}
\end{enumerate}
This will execute the precomputation generating 20 RDPF-triples, each of depth 20, with each instance working in three threads.

\begin{enumerate}
    \item \code{./prac -t 3 0 duoram 20 5}
    \item \code{./prac -t 3 1 localhost duoram 20 5}
    \item \code{./prac -t 3 2 localhost localhost duoram 20 5}
\end{enumerate}
Will execute the online phase, again with three threads each.
Specifying \code{duoram} will execute the tests as described in \ref{sub:testingAndBenchmarking} with 5 messages using RDPFs of depth 20 for each subtest.

\subsubsection*{General Execution}
All executions start in \code{prac.cpp} which accomodates the \code{main} method.
This includes the starting points for both the players' and the server's computations.
If parses the given arguments, sets the adresses of the other parties, and deceides whether the current execution is a player's or server's computation.
That means it calls \code{comp\_player\_main} or \code{server\_player\_main}.

\code{comp\_player\_main} represents the player's computation.
It is, again, separated into \code{preprocessing\_comp} and \code{online\_main} representing the preprocessing computation and the computations necessary for performing the onlien phase.

\code{server\_player\_main} represents the servers' computations.
It too is separated into its phases \code{preprocessing\_server} and \code{online\_main}.

The careful reader notices both the server and the player calling \code{comp\_main} albeit the previous text assumed these computations to be implemented separately.
Notice that this too is only a starting point for more specialized calculations.
Thus, \code{online\_main} too further specifies its computations given the number of the computational peer.

\subsubsection*{Preprocessing} \label{subsub:implementationPreprocessing}
%TODO wie wird das in files gespeichert und wo findet man die?


\subsubsection*{Cyclic Shifting} \label{subsub:implementationCyclicShifting}
\subsubsection*{Online Phase} \label{subsub:implementationOnlinePhase}
%TODO wie werden die DPFs (RDPF, CDPF) aufgerufen
\subsubsection*{Reading} \label{subsub:implementationReading}
\subsubsection*{Writing} \label{subsub:implementationWriting}

\section{General Limitations}
\label{sec:GeneralLimitations}
\subsection{AES and Label-Size}\label{sub:AESSizeLimitations}
The given code by Vadapalli et al. is optimized for 64-bit values with 128-bit labels.
These labels are represented by Intel's \code{\_\_m128i} SIMD-datatype (\ref{subsub:implementationAesni}).
Keys for AES are represented using an array holding eleven \code{\_\_m128i} values, which means they are of type \code{\_\_m128i[11]}.
Eleven values are necessary since AES for 128-bit values takes place in eleven rounds, thus eleven keys.
Adjusting AES for 256 or 512 bits requires more implementations using different inputs and keys, since a higher number of rounds is needed.

AES as used in the implementation keeps its secret-premises when concatenated.
This means, executing AES multiple times and concatenating the results into a single word is still considered secure.
Also, implementing arbitrarily scalable AES far exceeds this research's efforts.
Due to this reasons, the author chose to scale AES by concatenating 128-bit inputs.

This results in the requirements that all inputs, or the labels of given inputs, must be a multiple of 128 bits.
These inputs then can be splitted into 128-bit blocks, encrypted or decrypted using the 128-bit AES implementation, and finally concatenated to produce a single result.
This result is then the same length as the input, which is the equivalent of executing AES on the complete input instead of the divided block structure.

\subsection{WIDTH}\label{sub:WIDTH}
 %//TODO nochmal durch das folgende Kapitel durchgehen und prüfen ob sich Sachen wiederholen

\section{Possibilities}
\input{sections/implementation/possibilities/Possibilities.tex}

\subsection{Introducing a larger Register Type}
\label{sub:PossibilityLargerRegisterType}
Recall the register types \code{RegAS} and \code{RegXS} (\ref{itemize:importantStructs}).
They both represent a single register entry holding either an additive or XOR share of type \code{value\_t}.
The most trivial approach of representing values larger than \code{value\_t} can be acheived by combining register entries.
Therefore, this first approach considers new structs \code{BigAS}, \code{BigXS} implementing lists of registers.
Due to the need for different constructors, the author chose to introduce two new structs for index representation, namely \code{IndexAS} and \code{IndexXS}.

The author's adapted implementation is located on branch \code{NewInputType} on GitHub\footnote{The author's implementation of this approach is located here:\\ https://github.com/FelixWeik/PrivateRandomAccessComputations/tree/newInputType}.

\subsubsection*{BigAS, BigXS, and types.hpp} \label{subsub:BigASandBigXS}
These types represent a list of fixed length holding elements of type \code{RegAS} or \code{RegXS}.
Length restrictions are set in \code{types.hpp}.
It introduces new macros \code{INPUT\_BITS} and \code{INPUT\_PARTITION}.
\code{INPUT\_BITS} represents the bit-length of the value to represent.
It is set by the client and commits the execution to the chosen bitsize.
\code{INPUT\_PARTITION} represents the division $\frac{\text{INPUT\_BITS}}{\text{VALUE\_BITS}}$ which returns the length of the list.
Thus, by having \code{BigAS} or \code{BigXS} hold a list of \code{INPUT\_PARTITION} many elements, these lists reassemble a value of length \code{INPUT\_BITS}.
The general abstraction of \code{BigAS} and \code{BigXS} is called \textit{Large Register Type} in the following paragraphs.

A large register type is created using newly introduced constructors:
The standard constructor sets each list-element to 0. 
It can be populated anytime using the \code{set} method specifying the element and its position in the list.
Using a \code{RegAS} or \code{RegXS} pointer, respectively, populates the list using the copied values of the pointer.
Please ensure the amount of elements the pointer points to match the value of \code{INPUT\_PARTITION}.

Performing operations on either \code{BigAS} or \code{BigXS} must return results just like performing calculations on a vector of similar size.
Setting $\code{INPUT\_PARTITION} = 2$, for example, and multiplying it by $4$ must return the same result as performing the same calculation on a \code{uint128\_t} vector.
This marks the importance of adjusting the operators on these values.
This must also hold true for other operations like bit-shifts, additions, or logical operations.
Code snippets of \code{BigAS} are given in appendix \ref{lst:bigas}.

\subsubsection*{IndexAS and IndexXS} \label{subsub:IndexASandIndexXS}
Since the client wants to retrieve values of size \code{INPUT\_BITS}, the index type must adjust its size accordingly.
This will lead to her retrieving multiple entries of the database.
Each result is either of type \code{RegAS} or \code{RegXS}, wrapped in a list which itself is represented using a large register type.

Of course, one could use the same struct for \code{BigAS}, \code{IndexAS}, or \code{BigXS}, \code{IndexXS} respectively.
Two reasons led to the design decision of separating both structs.
First, the code becomes easier to read, test, and expand if the general structure is clear.
Next, indices and large register types require different constructors.
The default constructor for large register entries populates each entry with $0$.
This postpones the decision of what to write.
If an insufficient value of the wrong size is given, an error is thrown.
Indices must offer different behavior:
The default constructor populates the list with random values.
This means, \code{INPUT\_PARTITION} many random indices of the database are generated in this case.
Since the size of each index-element is of type \code{value\_t}, one cannot run into the error of reading from an index out of bounds.
Another index constructor takes a starting index as argument.
It generates $\text{INPUT\_PARTITION} - 1$ new indices beginning at the starting index.
This represents \code{INPUT\_PARTITION} many continuous reads or writes using this constructor.
An implementation example of \code{IndexAS} is given in appendix \ref{lst:indexas}.

\subsubsection*{Reading and Writing Operations using MemRefInd} \label{subsub:ReadingWritingUsingMemRefInd} 
As introduced in \ref{subsub:implementationCyclicShifting}, reading returns memory references, writing is executed on memory references.
Reading casts the reference into a register type, which is either \code{RegAS} or \code{RegXS}.
This is specified by using a template type \code{T} which is of either type.
For larger registers of type \code{BigAS} and \code{BigXS} reading and writing is especially performed on type \code{MemRefInd}.
It represents independent memory references which itself is a vector of explicit memory references.
These references can be accessed independently, reading and writing can be done in parallel.
Since the implemented larger register types and indices are not dependent themselves (unless they are part of a calculation), using independent memory references is the most performant option.
It also guarantees correctness and security as refreshing of the shape is automatically taken care of when using \code{MemRefInd} (\ref{subsub:largerRegisterTypeDPFNodesAndLeafLabels}).
The adapted operations using the new index types in context of reads and writes can be found in \ref{lst:bigDataAccess}.

\subsubsection*{DPF-Nodes and Leaf-Nodes}\label{subsub:largerRegisterTypeDPFNodesAndLeafLabels}
Recall the general concepts of \code{DPFNode} and \code{LeafNode} (\ref{itemize:importantAliases}).
A main question is whether \code{DPFNode}, initially set as \code{\_\_m128i}, needs to be scaled too when using different values for \code{INPUT\_PARTITION}.
Note that scaling \code{DPFNode} automatically scales \code{LeafNode} as well.

\ref{subsub:implementationDuoram} introduced the concepts of the \code{context} operation.
It is utilized by shapes for refreshing their context while keeping the underlying database untouched.
A single context is parametrized, for example, by the RDPF-Triple used for accessing a certain memory location.
This means updating the context uses a new set of RDPF-Triples.
Therefore, for every new operation on the database, this approach consumes a fresh set of RDPF-Triples.
Every element of the large index acts like a single access operation.
Thus, there is no need for scaling \code{DPFNode} and with that \code{LeafNode}.

\subsubsection*{Introducing new Tests} \label{subsub:IntroducingNewTests}
Test functions are implemented in \code{duoram.hpp}.
A new function \code{big()} executes independent reads and writes.
It is called using \code{big} instead of \code{duoram} while keeping the same notation of \ref{sub:implementationScenarios}.
Note that this test replaces solely the online phase, the preprocessing phase must be executed beforehand.
The notation for executing the preprocessing phase stays the same.
This approach is more RDPF hungry for reading and writing, consider generating more RDPF triples in the preprocessing phase.

More detailed information, tests, and benchmarking results is discussed further in \ref{sub:expEvalBenchmarks}.

\subsection{Increasing the Size of existing DPFs}
\label{sub:IncreasingTheSizeOfExistingDPFs}
Another approach is hardcoding larger datatypes or scaling existing datatypes replacing the originals.
The user chooses her preferred data-size before compilation and uses this specified size during execution.
This will commit each duoram to a single size which, however, can be chosen freely before compilation.

The updated version comes in two flavors:
\begin{enumerate}
    \item \code{64ToUnlimited}\footnote{Public on https://github.com/FelixWeik/PrivateRandomAccessComputations/tree/64ToUnlimited} as found on the \hyperlink{https://github.com/FelixWeik/PrivateRandomAccessComputations/tree/64ToUnlimited}{branch of the same name in the author's repository}.
    It uses a boilerplate-like approach of hardcoding the same functionality for multiple hardcoded datasizes.
    The \hyperlink{https://github.com/FelixWeik/PrivateRandomAccessComputations/tree/64ToUnlimited}{updated implementation} supports 64, 128, and 256 bit values.
    \item \code{UnlimitedGMP}\footnote{Public on https://github.com/FelixWeik/PrivateRandomAccessComputations/tree/UnlimitedGMP} utilizes the \textit{GNU Multi Precision Arithmetic Library (GMP)}\footnote{More details can be found here: https://de.wikipedia.org/wiki/GNU\_Multiple\_Precision\_Arithmetic\_Library}.
    In general, it allows for implementing integral datatypes like integers with arbitrary precision while supporting general bit-operations.
    Thus, the updated code found on branch \hyperlink{https://github.com/FelixWeik/PrivateRandomAccessComputations/tree/UnlimitedGMP}{UnlimitedGMP} supports any size being a multiple of 128 bits (\ref{sec:GeneralLimitations}).
    Further details are found in \ref{subsub:gmp} and \ref{subsub:applyingGMP}.
\end{enumerate}

Both approaches build upon the existing implementation and follow the same basic approach.
Yet, they differ significantly taking a deeper look into their codebase.

\subsubsection*{Locating necessary Changes} \label{subsub:necessaryChanges}
Considering the source code's structure, the author chose a top-down approach for implementing this solution.
This includes locating the most general datatypes, distinguishing cases and assigning each case its special datatype.
The most optimal solution could trigger a trickle-down-effect, where general datatypes affect more specialized datatypes with the ultimate goal of changing a single datatype or even a single flag in order to adjust the whole process of storing and retreiving data of a given size.
Therefore, the required efforts of achieving the optimal solution for this approach extend the efforts of solely switching the datatype used by the RDPFs.

The topmost elements in  this datatype-top-to-bottom structure are macros (\ref{subsub:implementationTypeshpp}).
These are specified in \code{types.hpp}.
They specify the most general datatypes and thus their derivates, aliases.
\begin{itemize}
    \item Type-alias \code{value\_t} by itself is used 258 times throughout the code.
    \code{VALUE\_BITS} calls the according cases.
    Its most prominent usage is in specifying the size of the register types \code{RegAS} and \code{RegXS}, and the size of DPF-leafs.
    \item \code{ADDRESS\_MAX\_BITS} specifies the size of addresses.
    The derived type-alias \code{address\_t} is used 143 times throughout the base-implementation.
    It deceides the type of indexing into the database.
    As the database is of size $2^\text{INPUT\_BITS}$ at least \code{INPUT\_BITS} many bits are needed for formatting adresses correctly.
\end{itemize}

\subsubsection*{Bitmasks and Bitutils}\label{subsub:bitmasksAndBitutils}
\code{bitutils.hpp} provides bitmasks and important functions for bit operations.
These operations include \code{xor\_if()} and \code{get\_lsb()} which are used multiple times throughout the code.
\code{xor\_if()} will execute a bitwise xor operation if a flag-bit is set, \code{get\_lsb()} retrieves the least significant bit of a given \code{\_\_m128i} value.

Bitmasks like \code{bool128\_mask} and \code{if128\_mask} are utilised when generating and traversing DPF-trees, during bit-operations, or while blinding shares.
In short, they increase bit-operations' performance by providing predefined values that will be used frequently during execution.

Since these bit-operations are executed primarily on \code{\_\_m128i} values, all function signatures and bitmasks are predefined for this specific datatype.
If one wants to scale their size arbitrarily, two solutions exist:
\begin{enumerate}
    \item Cut larger values into 128-bit blocks:
    This requires the input-type to be a multiple of 128-bit, which is not always the case.
    Also, it requires each function call to cut the input into $N$ blocks, call the given function $n$ times and finally concatenate each result to receive the final result.
    This too scales badly.
    \item Using GMP (\ref{sub:approachGMP}):
    By using GMP, one can create datatypes of arbitrary precision.
    Also, bit-operations are predefined for any \code{mpz} value.
    Thus, the specialized bit-operations defined in the code can be implemented to utilize the more general GMP bit-operations.
    Bitmasks can be left untouched, since the least significant bit or bit-inverses are not size-dependent.
    By using this library for bit-operations while leaving the bitmasks as they are, one achieves an arbitrarily scaled set of bit-utils.
    Therefore, the author further pursues this approach.
\end{enumerate}

The approach on GMP will be further discussed in the following section \ref{sub:approachGMP}.

\subsubsection*{Adjusting Compilation}\label{subsub:adjustingCompilation}
Both approaches require a flag bit to decide on which datatype or size to use during execution.
A new compilation-flag \code{-DVALUE\_BITS} is set in the \code{MAKEFILE} and represents the length of data one wants to store and retrieve.
It directly affects \code{VALUE\_BITS} and places the given value into it.
\code{VALUE\_BITS} then triggers the desired trickle-down affect which further changes macros and datatypes.
If nothing else is specified, \code{-DVALUE\_BITS} is set to $64$ by default which is similar to the base-implementation of Vadapalli et al.

\subsubsection*{Introducing new Tests}
%//TODO hier gehts weiter lieber Felix

\subsubsection*{General Discussion}

\subsection{DCF Intervals}
\label{sub:DCFIntervals}
Another possibility builds upon the already mentioned DCFs (\ref{sub:DistributedComparisonFunctions}).
Each read- and write-operation is expanded with 2 DCF-pairs.
These will be used to define a fixed-size interval.
Fixed-size for it is defined upon a fixed datasize.
This means, for example, a 128-bit value is stored in an interval of size two, 256-bit in an interval of size four and so on.

A high-level overview works as follows:
\begin{itemize}
    \item 
\end{itemize}

\section{Testing for Correctness}
\input{sections/implementation/TestingForCorrectness.tex}

\section{Performance Benchmarks}
%Anmerken, dass Referenzimplementierung aus Hardcodes für größere Datenwerte besteht